\documentclass{report}
\usepackage{amsmath,amssymb}
\setlength{\parindent}{0mm}
\setlength{\parskip}{1em}
\begin{document}
\begin{center}
\rule{6in}{1pt} \
{\large Abdulgader Almutairi \\
{\bf Enforcement of CA-UCON Model}}

Kingdom of Saudi Arabia \\ Alqassim ,Alrass
\\
{\tt Azmtierie@qu.edu.sa}\\
Saad Almutairi\end{center}

A Context-Aware Usage CONtrol (CA- UCON) model is an extension of the
traditional UCON model which enable adaptation to environmental changes
in the aim of preserving continuity of access. When the authorizations
and obligations requirements are met by the subject and the object, and
the conditions requirements fail due to changes in the environment or the
system con- text, CA-UCON model triggers specific actions to adapt to the
new situation. Besides the data protection, CA-UCON model so enhances the
quality of services, striving to keep explicit interactions with the user
at a minimum. In this paper, we propose an architecture of the reference
monitor for the CA-UCON model and investigate a variety of enforcement
approaches in ubiquitous computing systems; whether centralized,
distributed or hybrid; depending on applications.


\end{document}
